O uso de IA em diagnóstico oftalmológico tem avançado consideravelmente nas últimas décadas, sobretudo na detecção e classificação da retinopatia diabética. \cite{gulshan_development_2016} demonstraram que algoritmos baseados em redes neurais convolucionais podem detectar essa condição com precisão comparável à avaliação humana. \cite{abramoff_pivotal_2018} ampliaram esse uso, desenvolvendo sistemas capazes de quantificar o risco de progressão da doença, facilitando intervenções médicas personalizadas.

\cite{ting_development_2017} exploraram o potencial preditivo dessas tecnologias, identificando pacientes com risco elevado de progressão para formas avançadas da retinopatia. Esses avanços demonstram a eficácia e potencial escalabilidade da IA, especialmente no rastreamento populacional e em contextos de saúde pública, como já observado em países como a Inglaterra. Atualmente, inclusive, alguns sistemas baseados em IA como o IDx foram aprovados para uso clínico autônomo em cuidados primários \cite{topol_high-performance_2019}.

No Brasil, entretanto, ainda falta uma estratégia nacional padronizada, apesar da alta prevalência de RD e do impacto significativo sobre a população adulta, conforme destacado pelas diretrizes da Sociedade Brasileira de Diabetes (SBD). O experimento realizado em Sergipe mostrou desafios como baixa adesão mesmo com triagem gratuita, necessidade de adaptação tecnológica ao contexto local e ausência de marco regulatório específico para IA em saúde \cite{korn_malerbi_feasibility_2022}.